
\subsection{Soft Convex Hulls}
\label{sec:soft}
Let $W$ be a finite set of points in $\Re^d$, and
$T$ be the vertices of $\ch(W)$. The subset $T$ of $W$ is the unique subset of $W$ with the following properties: 
\begin{itemize}
\item[(P1)] $W\subseteq \ch(T)$
\item[(P2)] $\forall w\in W$, if $w\notin \ch(W\setminus \{w\})$, then $w \in T$. 
\end{itemize}



We now define a natural notion of {\em soft} convex hull. 
\begin{defn}\label{soft-CH}
For an $\varepsilon \geq 0$, and $S\subseteq W$, define the $\varepsilon$-convex hull of $S$, $\ech(S)$, as $\ch(S) + \varepsilon \dia(W) B$, where $B$ is the unit ball of the Eucliean norm. 
\end{defn}

 The intuition behind the above definition is that $\ch(W)$ can have a very large set of vertices, but there may be a small set of points whose soft convex hull contains $W$. This is defined more formally as follows: 

\begin{defn}\label{envelope}
We call a subset $T\subseteq W$ an $\varepsilon$-envelope of $W$, written
$\eenv(W)$, if 
$W\subseteq \ech(T).$  
\end{defn} 

\noindent
{\bf Remarks:} The following observations about the set $\eenv(W)$ are easy to see:
\begin{itemize}
\item[(a)] There are several distinct sets $T$ which could qualify as $\eenv(W)$. For example, let $W$ consist of the following set of points in $\Re^2$:  a set of points $W_1$ close to $(0,0)$ and a set of points $W_2$ close to $(1,0)$. Let $T$ be a pair of points $\{x,y\}$ with $x \in W_1, y \in W_2$. Then it is easy to check that $T$ is an $\eenv(W)$. 
\item[(b)] Let $T$ be $\eenv(W)$. Unlike property~(P2) above, it is not necessary that if $w \in W$ is such that $w \notin \ech(W \setminus \{w\})$, then $w \in T$. 
 \end{itemize}

Since the set $\eenv(W)$ is not uniquely determined, we will impose one more condition on it to make it unique (if it exists) and polynomial time computable. This condition requires
the points of $T$ to be ``far apart'' from each other. More precisely:

\begin{defn}\label{strict-envelope}
For $\varepsilon,\delta \in [0,1]$, a set $T$ is called a $\edenv(W)$ if it is an $\eenv(W)$ and  
\begin{equation}\label{792}
\forall w\in T, \dist(w,\ch(T\setminus \{w\}))> \delta \dia(W)
\end{equation}
\end{defn}

\begin{fact}
\label{f:check}
Given a subset $T$ of $W$, we can check in polynomial time whether $T$ is an $\edenv(W)$. 
\end{fact}

\begin{proof}
Fix a subset $T$. We can verify in polynomial time whether $T$ is an $\eenv(W)$. Indeed, for each point $w \in W$, we need to check if $w \in \ech(T)$. This can be expressed as a convex program feasibility problem, where the variables are $\lambda_t$ for each $t \in T$: 
$$ \left|\sum_{t \in T} \lambda_t \cdot t - w \right| \leq \varepsilon \dia(W), \quad \sum_t \lambda_t = 1, \quad \lambda_t \geq 0 \  \  \forall t \in T.$$ 

Similarly, we can check~\eqref{792} using convex programming. For each $w \in T$, we can find $\dst(w, \ch(T \setminus  \{w\})$ as follows, where the variables are $\lambda_x, x \in T \setminus \{w\}$: 
$$ \min. \ \left| \sum_{x \in T \setminus \{w\}} \lambda_x \cdot x - w \right|, \quad \sum_x \lambda_x = 1, \quad \lambda_x \geq 0 \  \  \forall x \in T\setminus \{w\}. $$
\end{proof}


For rest of the section, we address the following question: given the set $W$, and parameters $\varepsilon, \delta$, is there a  $\edenv(W)$, and if so, can we find in polynomial time an approximation to this set? We informally argue that several ``natural'' greedy strategies do not work. Consider, for example, the following algorithms for identifying $\edenv(W)$: 
\begin{itemize}
    \item Identify $T$ as the set of points $w\in W$ for which $w$ is not close to $\ch(W\setminus \{w\})$: We define a set of points in $\Re^2$. We have {\em rings} of points, one around $ (0,0)$ and the other around $(0,1)$, each point is close to the convex hull of the others and so this algorithm will   dismiss all the points as not belonging to the desired set $T$.
    % \item \alert{This example is not clear to me, i.e., I am not sure all the conditions on the 4 points can be satisfied. - rk   }Start by adding an extreme point (along a randomly chosen direction) to $T$. For the next $k-1$ steps, iteratively find a point $w$ for which $\dist(w,T)$ is maximized and add it to $T$: We consider a set of points in $\Re^2$ as follows: let $v_1, v_2, v_3$ be three affinely independent points with $|v_1 - v_2| > |v_1 - v_3|/\delta, |v_2 - v_3|/\delta$ and $\dist(v_3, \ch(\{v_1, v_2\})) > \delta |v_1 - v_2|$. Now let $v_4$ be the mid-point of the line segment joining $(v_1, v_2)$. Let $W=\{v_1, v_2, v_3, v_4\}$. Say the first point selected by the greedy algorithm is $v_1$. Then the next two points selected by it will be $v_2, v_4$ respectively. But the ``correct'' $T$ is $\{v_1, v_2, v_3\}$. 
    \item Start by adding an extreme point (along an arbitrarily chosen direction) to $T$. For the next $k-1$ steps, iteratively find a point $w$ for which $\dist(w,\ch(T))$ is maximized and add it to $T$: a simple example as above shows that this idea may not work as well. Consider $v_1, v_2, v_3, v_4$ lying on corners of a square, and let $v_5$ be the mid-point of $v_3, v_4$. We start by adding $v_1$ and then $v_2$ to $T$. But in the next step, we could add $v_5$, and then both $v_3, v_4$. But the ``correct'' $T$ would have been $\{v_1, v_2, v_3, v_4\}$. 
\end{itemize}
 


If $\varepsilon=0$, the above question is easy to answer in polynomial time. The answer is yes iff the set $T$ of vertices of $\ch(W)$ satisfies (\ref{792}). Also, if $\delta=1$, $T$ has to be a singleton to satisfy (\ref{792}). 

In rest of this section, we consider the following problenm: For what pairs of values of $\varepsilon,\delta,$ can we prove that there is {\em essentially} at most one $\edenv(W)$, and if so, can we determine this set efficiently ? 
We do not know the exact answer to this, but our main result here (which suffices for the applications) is 
(verbally stated) an affirmative answer to the question if the following condition is satisfied:
$$\delta\in \Omega(\sqrt\varepsilon).$$
This will follow as a corollary of our main result:

\begin{theorem}
\label{thm:soft}
Let
$\delta,\varepsilon,\varepsilon_3$ be reals in $(0,1/8)$ satisfying
\begin{equation}
\label{822}\delta> \max \left(\frac{2\varepsilon}{\varepsilon_3-\varepsilon}, 4 \varepsilon_3 \right)
\end{equation}
Let $W$ be a finite set of points in $\Re^d$. 
We can 
determine in polynomial time whether  exists a set $T$ in $\edenv(W)$, and if so,  we can efficiently 
find a subset $Q$ of $W$ such that
\begin{align}
|Q|=|T|\\
\forall w\in T, \exists x\in Q: |w-x|\leq 2\varepsilon_3 \dia(W)
\end{align}
\end{theorem}
\begin{corollary}\label{cor:soft}
Let
$\delta,\varepsilon$ be reals in $(0,1/8)$ satisfying $\delta>16\sqrt\varepsilon$
Let $W$ be a finite set of points in $\Re^d$. 
We can 
determine in polynomial time whether  exists a set $T$ forming a $\edenv(W)$, and if so,  we can efficiently 
find a subset $Q$ of $W$ such that
\begin{align}
|Q|=|T|\\
\forall w\in T, \exists x\in Q: |w-x|\leq 8\sqrt \varepsilon \dia(W)
\end{align} 
\end{corollary}
\begin{proof}
The Corollary follows from Theorem (\ref{thm:soft}) by taking $\varepsilon_3=4\sqrt\varepsilon$ 
\end{proof}
\begin{proof} (of \Cref{thm:soft})
The procedure is described  in~\Cref{algo:main}. We first compute a subset $Q''$ of $W$ consisting of points $w \in W$ which do not lie in the soft convex hull of the points in $W$ which are ``far'' from $w$ -- this can be doen in polynomial time by using arguments similar to those in the proof of~\Cref{f:check}. Then $Q$ is defined as a maxumal subset of points in $Q''$ with the property that the pair-wise distance between the points in it large. Finally, we check whether $Q$ is an $\edenv(W)$, which can be done efficiently using~\Cref{f:check}. 


\begin{algorithm}[H]
  \caption{Procedure to identify $\edenv(W)$}
  \label{algo:main}
    Compute $\dia(W)$. \;
    Let $Q' := \{w \in W: w \in \ech \left(\{x \in W: |w-x| \geq \varepsilon_3 \dia(W) \}\right). $ \;
    Let $Q'' := W \setminus Q'$. \;
    Define $Q :=$ maximal subset of $Q''$ such that for every distinct $x,y \in Q$, $|x-y| > 2\varepsilon_3 \dia(W). $ \;
    \If{$Q$ is $\edenv(W)$} {
       Output  the set $Q$. \; }
    \Else {
     Output No. }
\end{algorithm}


Now we analyze this algorithm. If there is no $\edenv(W)$, then the algorithm will clearly say ``No'' (because the set $Q$ will not be $\edenv(W)$). So assume that there is a set $T$ which is $\edenv(W)$. 

\begin{claim}
\label{cl:ed}
$T \subseteq Q''$. 
\end{claim}
\begin{proof}

Suppose for the sake of contradiction that there is a point $w \in T \setminus Q''$.  For sake of brevity, let $W_x$ denote the subset of points in $W$ satisfying $|w-x|  \geq \varepsilon_3 \dia(W)$. The fact that $w \notin Q''$ implies that $w \in \ch(W_x)$, i.e., 
$$ w = \sum_{x \in W_x} \lambda_x \cdot x + e_0, $$
where $\lambda_x$ form a convex combination and $|e_0| \leq \varepsilon \dia(W).$ Since $W \subseteq \eenv(T)$, for each $x \in W_x$, there are a point $x' \in \ch(T)$ such that $e_x := x-x'$ has length at most $\varepsilon \dia(W)$. Since $w \in T$, we can write $x'$ as 
$$ x' = \mu_x w + (1-\mu_x) y_x, \quad y_x \in \ch(T \setminus \{w\}), $$
where $\mu_x \in [0,1]$. 
Therefore, 
$$ |w-x| \leq |w-x'| + |e_x| = (1-\mu_x) |w-y_x| + |e_x| \leq (1-\mu_x) \dia(W) + \varepsilon \dia(W). $$
But we know that $|w-x| \geq \varepsilon_3 \dia(W)$. Therefore, we get
\begin{align}
\label{eq:111}
(1- \mu_x) \geq \varepsilon - \varepsilon_3. 
\end{align}
Now, 
\begin{align*}
w & = \sum_{x \in W_x} \lambda_x \cdot x + e_0 = \sum_{x \in W_x} \lambda_x \cdot x' + \sum_x \lambda_x e_x + e_0 \\
& = \sum_{x \in W_x} \lambda_x \mu_x \cdot w + \sum_{x \in W_x} \lambda_x (1 - \mu_x) y_x + e,  
\end{align*}
where $e = \sum_x \lambda_x e_x + e_0$ has length at most $2\varepsilon \dia(W)$. 
Let $\theta_x$ denote $\lambda_x(1-\mu_x)$. Observe that $\theta:= \sum_x \theta_x = 1 - \sum_x \lambda_x \mu_x$. Therefore, we can rewrite the above as 
$$ w - \frac{1}{\theta} \sum_{x \in W_x} \theta_x y_x = \frac{e}{\theta}. $$
Since $\frac{1}{\theta} \sum_{x \in W_x} \theta_x y_x \in \ch(T \setminus \{w\})$, we see that 
$$ \dist(w, \ch(T \setminus \{w\}) \leq \frac{|e|}{\theta} \leq 
\frac{2\varepsilon \dia(W)}{\sum_{x \in W_x} \lambda_x (1-\mu_x)} 
\stackrel{\eqref{eq:111}}{\leq} \frac{2\varepsilon}{\varepsilon-\varepsilon_3} \dia(W). $$
Using~\eqref{792} and~\eqref{822}, we get a contradiction. Therefore, $w \in Q''$. 
\end{proof}

Since $T \subseteq Q''$, for every $w \in T$, there exists a point $x_w \in Q$ such that $|w-x_w| \leq 2 \varepsilon_3 \dia(W). $ It is easy to see that if $w, w'$ are two distinct points in $T$, then $x_w \neq x_{w'}$. Indeed, 
$$ |x_w - x_{w'}| \geq |w-w'| - |w-x_w| - |w'-x_{w'}| \geq \delta \dia(W) - 4 \varepsilon_3 \dia(W) \stackrel{\eqref{822}}{>} 0. $$
Finally, we prove that every element in $Q$ is of the form $x_w$ for some $w \in T$. 

\begin{claim}
\label{cl:final}
Every $x \in Q$ is of the form $x_w$ for some $w \in T$. 
\end{claim}
\begin{proof}
Consider a point $x \in Q$. This implies that $x \in Q''$. We claim that there is a point $w \in T$ such that $|x-w| \leq \varepsilon_3 \dia(W)$. Suppose not. Then the set of points $w$ such that $|x-w| \geq \varepsilon_3 \dia(W)$ includes $T$. But we know that $x \in \ech(T)$, and so $x \in Q'$, a contradiction. 

Therefore $$|x - w_x| \leq |x- w| + |w-w_x| \leq  2\varepsilon \dia(W). $$
This implies that $x=w_x$ (since any two distinct points in $Q$ have distance greater than $2\varepsilon \dia(W)$. 
\end{proof}

Thus we have shown that $|Q|=|T|$ and for every $w \in T$, there is a unique element $x_w \in Q$ with $|w-x_w| \leq 2 \varepsilon_3 \dia(W)$.
\end{proof}


